\section{Keterkaitan Buku Ajar dan RPS OBE}

Buku ajar ini dirancang dengan ketat mengikuti Rencana Pembelajaran Semester (RPS) Mata Kuliah Kriptografi yang berbasis \textbf{Outcome-Based Education (OBE)}. Setiap bab dipetakan secara eksplisit ke Capaian Pembelajaran Mata Kuliah (CPMK) dan Sub-CPMK yang tercantum dalam silabus.

Outcome-Based Education (OBE) merupakan pendekatan kurikulum yang berfokus pada hasil pembelajaran (learning outcomes) yang harus dicapai mahasiswa, bukan sekadar materi yang diajarkan oleh dosen. Menurut Wikipedia, OBE mengorganisir seluruh sistem pendidikan menuju apa yang dianggap esensial bagi pembelajar untuk berhasil dilakukan di akhir pengalaman belajar mereka. Pendekatan ini memberikan kejelasan tentang harapan pencapaian dan fleksibilitas dalam metode instruksional, memungkinkan dosen untuk menggunakan berbagai teknik pengajaran dan asesmen yang sesuai dengan kebutuhan mahasiswa.

Implementasi OBE dalam kriptografi memastikan setiap kompetensi yang dipelajari dapat diukur dan diverifikasi melalui asesmen yang terstruktur. Model pembelajaran berbasis outcome ini memungkinkan perbandingan standar kompetensi antar institusi dan memfasilitasi transisi mahasiswa ke dunia kerja dengan bukti pencapaian kompetensi yang jelas. Mahasiswa menjadi lebih terlibat dan bertanggung jawab atas pembelajaran mereka sementara dosen berperan sebagai fasilitator yang mengarahkan pencapaian kompetensi yang telah ditetapkan.

\begin{table}[htbp]
\centering
\small
\begin{tabular}{|c|c|p{7cm}|}
\hline
\textbf{Bab} & \textbf{Sub-CPMK} & \textbf{Materi Terkait} \\
\hline
I & -- & Orientasi dan panduan belajar \\
\hline
II & 1.1, 1.2, 1.3 & Landasan teori, model keamanan, simetris vs asimetris \\
\hline
III & 2.1, 2.2 & Matematika kriptografi, Euclid, fast exponentiation \\
\hline
IV & 3.1, 3.2 & Caesar, Vigenère, substitusi, implementasi C \\
\hline
V & 4.1 & DES, AES, mode operasi block cipher \\
\hline
VI & 4.2 & RSA, implementasi C \\
\hline
VII & 5.1, 5.2 & Fungsi hash, tanda tangan digital \\
\hline
VIII & 6.1 & Protokol Diffie-Hellman, manajemen kunci \\
\hline
IX & 7.1 & Kriptanalisis, frequency analysis, Kasiski \\
\hline
X & 8.1, semua & Evaluasi, refleksi, integrasi kompetensi \\
\hline
\end{tabular}
\caption{Petakan Bab ke Sub-CPMK}
\end{table}

Pendekatan OBE menekankan bahwa setiap aktivitas pembelajaran---materi, latihan, praktikum, dan asesmen---harus diarahkan untuk mencapai kompetensi yang terukur. Dengan demikian, mahasiswa dapat memantau kemajuan mereka secara mandiri melalui checklist pencapaian yang disediakan di setiap bab.
